Every figure the pipeline emits, documented by scope, reading procedure, the data-availability problems that distort it on sparse runs, and the justification for its chart type --- grouped by metric family.

The pipeline emits up to 21 figures, registered by name in \code{PLOT\_DESCRIPTIONS} and written by \code{maybe\_make\_plots}. This catalogue documents each one along the four axes that decide whether a figure earns its place: its \textbf{scope}(what slice of data it shows),\textbf{how to read it}, the \textbf{data-availability problems} that distort it on sparse runs, and the \textbf{justification} for that chart type over the alternatives.

\begin{bmkeybox}{One reading rule that governs half the figures}
Several plots are bounded [0,1] and zero-inflated: on weak models the bars genuinely collapse to zero, and a zero can mean either "good" (no hallucination) or "nothing produced" (no actions emitted). Whenever a bar is flat-zero, disambiguate against SR and plan counts before reading it as a strength. The atlas notes this per-plot where it can cause problems.
\end{bmkeybox}

\subsubsection{Grounding / discriminatory}

\begin{bmplotcard}{Action hallucination (bars)}
\code{hallucination\_by\_model\_domain}
\begin{description}
\item[chart] grouped bars
\item[scope] Same action-hallucination rate, grouped bars per model across domains.
\item[how to read] Bar height = fraction of emitted actions whose name is not in the domain schema.
\item[data problems] Zero-height bars are ambiguous: genuinely grounded vs nothing emitted. Cross-check plan counts.
\item[why this form] Bars expose per-domain spread for one model better than a heatmap row when domain count is small.
\end{description}
\end{bmplotcard}

\begin{bmplotcard}{Object hallucination}
\code{object\_hallucination\_by\_model\_domain}
\begin{description}
\item[chart] grouped bars
\item[scope] Fraction of action arguments referencing objects absent from the problem file.
\item[how to read] Separates' wrong object' errors from' wrong action' errors; compare against the action-hallucination twin.
\item[data problems] Undefined when a model emits no arguments; appears as 0 and must not be read as' good'.
\item[why this form] Splitting object- from action-hallucination is the whole point --- two distinct grounding failures need two plots.
\end{description}
\end{bmplotcard}

\subsubsection{Execution \& reasoning}

\begin{bmplotcard}{Executability ratio}
\code{executability\_by\_model\_domain}
\begin{description}
\item[chart] box / strip
\item[scope] Distribution of exec\_ratio (valid-prefix length / plan length) per model\(\times\)domain.
\item[how to read] Higher = plans run further before the first precondition break. 1.0 = fully executable (often also valid).
\item[data problems] Single-instance domains give a degenerate' distribution' of one point; spread is only meaningful with several instances.
\item[why this form] A distribution view (not a mean) is justified because exec\_ratio varies instance-to-instance and the spread is informative.
\end{description}
\end{bmplotcard}

\begin{bmplotcard}{Executability vs plan length}
\code{executability\_vs\_length}
\begin{description}
\item[chart] scatter (faceted)
\item[scope] exec\_ratio against plan length, one facet per domain.
\item[how to read] A downward trend = the model breaks down on longer horizons (the Embers-of-Autoregression length cliff).
\item[data problems] Needs many instances spanning a length range; sparse runs show a near-empty scatter with no visible trend.
\item[why this form] Scatter is the only honest way to show a per-instance relationship between two continuous quantities.
\end{description}
\end{bmplotcard}

\begin{bmplotcard}{Failure type breakdown}
\code{failure\_type\_breakdown}
\begin{description}
\item[chart] stacked bars
\item[scope] Per model: count of sequencing errors vs state fabrications among failed plans.
\item[how to read] Tall sequencing segment = grounded model failing on order; tall fabrication segment = invents unsatisfiable states.
\item[data problems] Models with almost no failures (or no plans) give tiny/empty bars --- proportion is noisy at low counts.
\item[why this form] Stacking the two mutually exclusive failure causes shows both the mix and the absolute volume in one bar.
\end{description}
\end{bmplotcard}

\begin{bmplotcard}{Temporal distance}
\code{temporal\_distance\_by\_model}
\begin{description}
\item[chart] bars
\item[scope] Mean gap (steps) between the last true-precondition step and the actual failure step, sequencing errors only.
\item[how to read] Larger = the model stays internally consistent for longer before diverging (a' compositional reach' proxy).
\item[data problems] Defined only for sequencing errors; models with none get NaN/absent bars --- absence is not zero.
\item[why this form] A simple per-model bar suffices because temporal distance is already a per-model mean of a single quantity.
\end{description}
\end{bmplotcard}

\subsubsection{Planning efficiency}

\begin{bmplotcard}{FASR by model\(\times\)domain}
\code{fasr\_by\_model\_domain}
\begin{description}
\item[chart] grouped bars
\item[scope] First-attempt success rate per model across domains --- success with no repair.
\item[how to read] The retry-free skill signal; compare against SR bars to see how much each model leans on retries.
\item[data problems] Bounded [0,1]; zero bars common and meaningful (no first-try solves).
\item[why this form] Bars per domain make the zero-inflation and the few non-zero domains immediately legible.
\end{description}
\end{bmplotcard}

\begin{bmplotcard}{FASR by difficulty}
\code{fasr\_by\_difficulty}
\begin{description}
\item[chart] bars
\item[scope] First-attempt success rate aggregated by difficulty tier.
\item[how to read] A monotone decline across tiers is the expected difficulty gradient; a cliff signals a capability boundary.
\item[data problems] Requires a populated difficulty field; if tiers are unlabelled or unbalanced the bars mislead.
\item[why this form] Binning by difficulty is the natural way to expose the McCoy difficulty cliff.
\end{description}
\end{bmplotcard}

\begin{bmplotcard}{IWSR by model\(\times\)domain}
\code{iwsr\_by\_model\_domain}
\begin{description}
\item[chart] grouped bars
\item[scope] Iteration-weighted success rate (1/iterations on success, else 0) per model\(\times\)domain.
\item[how to read] Discounts late successes; sits between FASR and SR. Closeness to SR = efficient repair.
\item[data problems] Bounded [0,1]; like FASR it is zero-inflated for weak models.
\item[why this form] Same grid logic as FASR; bars keep the three efficiency metrics visually comparable.
\end{description}
\end{bmplotcard}

\begin{bmplotcard}{SR \(\cdot\) FASR \(\cdot\) IWSR triptych}
\code{sr\_fasr\_iwsr\_by\_model}
\begin{description}
\item[chart] grouped bars
\item[scope] All three success metrics side by side per model.
\item[how to read] Read the descent SR\(\rightarrow\)IWSR\(\rightarrow\)FASR: a steep drop = retry-inflated; a flat triple = genuine first-try skill.
\item[data problems] For SR\(\approx\)0 models all three vanish together, conveying little beyond' failed'.
\item[why this form] Co-plotting the three is the cleanest single picture of retry dependence per model.
\end{description}
\end{bmplotcard}

\begin{bmplotcard}{Retry gap}
\code{retry\_gap\_by\_model}
\begin{description}
\item[chart] bars
\item[scope] SR - FASR per model; the Stechly retry-dependence signal.
\item[how to read] Higher = more of the reported success came only after retries.
\item[data problems] Zero gap is ambiguous: either pure first-try skill or total failure (SR=FASR=0). Disambiguate with SR.
\item[why this form] A single derived bar per model is appropriate; the gap is one number with a clear sign.
\end{description}
\end{bmplotcard}

\begin{bmplotcard}{P(Valid | reached k)}
\code{p\_valid\_given\_k}
\begin{description}
\item[chart] line
\item[scope] Probability a plan is valid given the model reached attempt k, across k.
\item[how to read] Rising curve = Efficient Corrector (later attempts really are better); flat line = Stochastic Searcher (resampling).
\item[data problems] Late-k points rest on few surviving instances --- the right tail is high-variance and should be read cautiously.
\item[why this form] A line over k is the only form that reveals the shape (rising vs flat) the profiles depend on.
\end{description}
\end{bmplotcard}

\subsubsection{CoT alignment}

\begin{bmplotcard}{CoT alignment}
\code{cot\_alignment\_by\_model\_domain}
\begin{description}
\item[chart] grouped bars
\item[scope] Mean CoTaL (reasoning-vs-plan alignment) per model\(\times\)domain.
\item[how to read] Higher = the reasoning trace structurally matches the emitted plan.
\item[data problems] Mostly NaN: most instances have no separate reasoning plan, so bars are absent or rest on tiny samples.
\item[why this form] Bars are fine where data exists; the dominant story here is missingness, documented in the CoT view.
\end{description}
\end{bmplotcard}

\subsubsection{Composite \& cross-domain}

\begin{bmplotcard}{Success-rate heatmap}
\code{success\_rate\_heatmap}
\begin{description}
\item[chart] heatmap
\item[scope] Raw SR per model\(\times\)domain --- no normalisation, so cells are comparable across the whole grid.
\item[how to read] Empty columns = domains nobody solved (difficulty or infrastructure); empty rows = models that solved nothing.
\item[data problems] Genuinely sparse: many cells are 0, which is informative rather than missing.
\item[why this form] Unlike the rank heatmap, raw SR is on one common scale, so a single-scale heatmap is valid.
\end{description}
\end{bmplotcard}

\begin{bmplotcard}{Composite Planning Score}
\code{composite\_scores}
\begin{description}
\item[chart] bars + CI
\item[scope] PS per model with the 95\% bootstrap CI.
\item[how to read] Sort by height for the leaderboard; overlapping CIs = the difference is not resolvable.
\item[data problems] Overall PS can sit outside its own per-row bootstrap CI for NaN-heavy models (kimi caveat).
\item[why this form] Bars with error bars are the standard, honest way to show a ranked scalar with uncertainty.
\end{description}
\end{bmplotcard}

\begin{bmplotcard}{Within-domain rank heatmap}
\code{domain\_ranking\_heatmap}
\begin{description}
\item[chart] heatmap (faceted)
\item[scope] Per-domain normalised rank (0=best,1=worst) for every RANK\_METRIC.
\item[how to read] Compare cells DOWN a column WITHIN one domain block only --- never across domain blocks.
\item[data problems] Normalisation is per-domain, so identical colours in different blocks are not comparable scores.
\item[why this form] Per-domain normalisation onto one colour scale is what lets many domains share a figure.
\end{description}
\end{bmplotcard}

\begin{bmplotcard}{Rank variance}
\code{rank\_variance}
\begin{description}
\item[chart] bars
\item[scope] Mean std of within-domain rank across domains, per model.
\item[how to read] Higher = domain specialist; lower = generalist.
\item[data problems] Computed on raw ranks (not normalised) --- magnitudes are run-specific; ordering is robust.
\item[why this form] A ranked bar is the right form for a single per-model dispersion scalar.
\end{description}
\end{bmplotcard}

\begin{bmplotcard}{Domain correlation (Spearman)}
\code{domain\_correlation}
\begin{description}
\item[chart] heatmap
\item[scope] Model\(\times\)model Spearman \(\rho\) of per-domain success-rate vectors.
\item[how to read] \(\rho\)\(\approx\)1 = redundant probes; low/negative \(\rho\) = complementary models worth keeping both.
\item[data problems] Needs several domains with variance; with few domains and many SR=0 ties \(\rho\) is dominated by ties.
\item[why this form] A correlation matrix is the canonical view for pairwise similarity across a model set.
\end{description}
\end{bmplotcard}

\begin{bmplotcard}{PS stacked by domain}
\code{ps\_by\_domain\_stacked}
\begin{description}
\item[chart] stacked bars
\item[scope] Each model's PS split into per-domain contributions.
\item[how to read] Same total height can be one tall segment (concentrated) or many even ones (broad) --- read the segmentation.
\item[data problems] Domains with no PS contribution simply don't appear; ensure the domain set is comparable across models.
\item[why this form] Stacking is the natural encoding for "is this total broad or concentrated" --- a pie would hide the total.
\end{description}
\end{bmplotcard}

\begin{bmplotcard}{Metrics summary table}
\code{metrics\_summary\_table}
\begin{description}
\item[chart] table figure
\item[scope] All key aggregate metrics rendered as a matplotlib table for archiving.
\item[how to read] A static snapshot of the overall table; read it as the numeric source of the bar/heatmap figures.
\item[data problems] None beyond whatever NaNs the underlying tables carry; it faithfully shows them.
\item[why this form] A rendered table travels with the plot bundle so a figure folder is self-documenting.
\end{description}
\end{bmplotcard}

\begin{bmplotcard}{Capability radar}
\code{radar\_chart}
\begin{description}
\item[chart] radar
\item[scope] Polar pentagon of FASR \(\cdot\) IWSR \(\cdot\) Exec \(\cdot\) IHR \(\cdot\) PAS per model (base profile, no CoT bonus).
\item[how to read] Area = overall strength; shape = capability signature. A collapsed pentagon = No Grounding.
\item[data problems] Equal areas can hide very different shapes; many overlaid models become unreadable --- facet for publication.
\item[why this form] A radar is justified for a small fixed set of bounded axes where the shape itself is the message.
\end{description}
\end{bmplotcard}

\begin{bmcavbox}{Why the PNGs are not embedded here}
The 21 figures are run-specific --- they depend on which models and domains were in the bundle and would be stale the moment the suite is re-run. This catalogue documents what each figure \emph{means and how to read it} so it stays valid across runs; regenerate the actual PNGs with the pipeline against your own \code{results/plots/} directory.
\end{bmcavbox}
